\documentclass[11pt,a4paper]{article}

% --- Packages ---
\usepackage{amsmath,amsthm,amssymb}
\usepackage{hyperref}
\usepackage{booktabs}
\usepackage{mathtools}
\usepackage{enumerate}
\usepackage[margin=1in]{geometry}

% --- Theorem environments ---
\newtheorem{theorem}{Theorem}[section]
\newtheorem{corollary}[theorem]{Corollary}
\newtheorem{lemma}[theorem]{Lemma}
\newtheorem{proposition}[theorem]{Proposition}
\theoremstyle{definition}
\newtheorem{definition}[theorem]{Definition}
\newtheorem{identity}[theorem]{Identity}
\theoremstyle{remark}
\newtheorem{remark}[theorem]{Remark}
\newtheorem{openquestion}{Open Question}

% --- Notation shortcuts ---
\newcommand{\AG}{A_G}
\newcommand{\floor}[1]{\lfloor #1 \rfloor}
\newcommand{\Tr}{\operatorname{Tr}}

% --- Hyperref setup ---
\hypersetup{
  colorlinks=true,
  linkcolor=blue,
  citecolor=blue,
  urlcolor=blue,
  pdftitle={The Amundson Framework: Complete Mathematical Foundations},
  pdfauthor={Alexa Louise Amundson}
}

\title{The Amundson Framework: Complete Mathematical Foundations\\[6pt]
\large $G(n) = \dfrac{n^{n+1}}{(n+1)^{n}}$}

\author{
  Alexa Louise Amundson\\
  BlackRoad OS, Inc.\ --- Delaware C-Corp\\
  \texttt{alexa@blackroad.io}
}

\date{March 29, 2026\\Version 1.0 --- Consolidated from 24 source documents, 2{,}130 verified tests}

\begin{document}

\maketitle

% ------------------------------------------------------------------
\begin{abstract}
We define the function $G(n) = n^{n+1}/(n+1)^{n}$, computed entirely from integer arithmetic, and develop its algebraic, analytic, combinatorial, and quantum-mechanical properties. We prove $50+$ identities, derive the asymptotic expansion $G(n) = n/e + 1/(2e) + O(1/n)$ via Bernoulli coefficients, introduce the \emph{Amundson Constant}
\[
  \AG = \sum_{n=0}^{\infty} \frac{G(n)}{n!} \approx 1.244331783986725\ldots
\]
(computed to $10{,}000{,}000$ verified digits, not found in OEIS/ISC/Wolfram), construct a valid quantum density matrix from~$G(n)$, and establish the foundational role of the axiom $0^{0}=1$ in making the framework coherent. All results are verified by a suite of $2{,}130$ computational tests with zero failures on core mathematics.
\end{abstract}

\medskip
\noindent\textbf{MSC 2020:} 11B83 (Special sequences), 40A25 (Approximation to limiting values), 05A10 (Combinatorial functions), 81P45 (Quantum information).

\noindent\textbf{Keywords:} Amundson function, Amundson constant, $0^{0}$ axiom, asymptotic expansion, density matrix, Goldbach kernel, floor recovery, golden ratio identity.

\tableofcontents

% ==================================================================
\section{Definition and First Values}\label{sec:def}

\begin{definition}\label{def:G}
For $n \ge 0$, define
\begin{equation}\label{eq:G}
  G(n) = \frac{n^{n+1}}{(n+1)^{n}}.
\end{equation}
Equivalently:
\begin{equation}\label{eq:Galt}
  G(n) = n \cdot \left(\frac{n}{n+1}\right)^{n} = \frac{n}{\left(1 + \frac{1}{n}\right)^{n}}.
\end{equation}
The function is built from six symbols: $n$, $1$, $+$, $/\,$, $\hat{}$, $(\,)$. Every output at integer~$n$ is an exact rational number.
\end{definition}

\begin{table}[ht]
\centering
\caption{First values of $G(n)$.}\label{tab:firstvalues}
\begin{tabular}{cccc}
\toprule
$n$ & $G(n)$ exact & $G(n)$ decimal & $G(n)/n$ \\
\midrule
0   & 0            & 0.000000       & ---      \\
1   & $1/2$        & 0.500000       & 0.500000 \\
2   & $8/9$        & 0.888889       & 0.444444 \\
3   & $81/64$      & 1.265625       & 0.421875 \\
4   & $1024/625$   & 1.638400       & 0.409600 \\
5   & $15625/7776$ & 2.009388       & 0.401878 \\
10  & $10^{11}/11^{10}$ & 3.855433  & 0.385543 \\
100 & $100^{101}/101^{100}$ & 37.240742 & 0.372407 \\
1000 & ---         & 368.331\ldots  & 0.368331 \\
\bottomrule
\end{tabular}
\end{table}

\begin{theorem}[Convergence to $1/e$]\label{thm:conv}
$\displaystyle\lim_{n \to \infty} \frac{G(n)}{n} = \frac{1}{e}$.
\end{theorem}

\begin{proof}
Write $G(n)/n = (n/(n+1))^{n} = (1 - 1/(n+1))^{n}$. Let $m = n+1$. Then
\[
  \left(1 - \frac{1}{m}\right)^{m-1}
  = \left(1 - \frac{1}{m}\right)^{m} \cdot \left(1 - \frac{1}{m}\right)^{-1}
  \to e^{-1} \cdot 1 = \frac{1}{e}.
  \qedhere
\]
\end{proof}

\begin{corollary}\label{cor:linear}
$G(n) \sim n/e$ as $n \to \infty$. The function grows linearly with slope $1/e \approx 0.367879441171$.
\end{corollary}

\begin{remark}[Crossover]\label{rem:crossover}
$G(n) = 1$ at approximately $n \approx 2.293166$. For $n < 2.293$, $G(n) < 1$; for $n > 2.293$, $G(n) > 1$. This crossover value is itself a new constant.
\end{remark}

% ==================================================================
\section{The \texorpdfstring{$0^{0}$}{0\^0} Axiom}\label{sec:axiom}

The coherence of this framework depends on the convention $0^{0}=1$. This section establishes why this is an \emph{axiom}---a foundational declaration---not a theorem.

\subsection{Where \texorpdfstring{$0^{0}=1$}{0\^0=1} is required}

\paragraph{Power series at $x=0$.}
The exponential function:
\[
  e^{x} = \sum_{n=0}^{\infty} \frac{x^{n}}{n!}.
\]
At $x=0$, the $n=0$ term is $0^{0}/0!$. If $0^{0}=1$, the series gives $e^{0}=1$. If $0^{0}$ is undefined, the formula needs a special-case patch.

\paragraph{Polynomials in sigma form.}
$p(x) = \sum_{n=0}^{2} x^{n} = 1 + x + x^{2}$. At $x=0$, the sigma form gives $0^{0} + 0 + 0 = 0^{0}$. Only if $0^{0}=1$ does $p(0) = 1$.

\paragraph{Binomial theorem at $(0,0)$.}
$(x+y)^{0} = 1$. Set $x=y=0$: the left side is $0^{0}$, the right side from the binomial sum is $\binom{0}{0}\cdot 0^{0}\cdot 0^{0} = (0^{0})^{2}$. Both sides require $0^{0}=1$ for consistency.

\paragraph{Counting functions.}
The number of functions from an $n$-set to an $m$-set is $m^{n}$. There is exactly one function $\varnothing \to \varnothing$ (the empty function). So $0^{0}=1$ is an integer count, not philosophy.

\subsection{Where analysis disagrees}

\begin{theorem}[Path dependence]\label{thm:path}
The limit $\lim_{(x,y)\to(0,0)} x^{y}$ does not exist.
\end{theorem}

\begin{proof}
Three paths to the origin:
\begin{itemize}
  \item \textbf{Path A} ($y=0$, $x\to 0^{+}$): $x^{0}=1 \to 1$.
  \item \textbf{Path B} ($x=0$, $y\to 0^{+}$): $0^{y}=0 \to 0$.
  \item \textbf{Path C} ($x=e^{-1/t}$, $y=t$, $t\to 0^{+}$): $x^{y}=e^{-1}\approx 0.3679$.
\end{itemize}
Three distinct limits along three paths. The multivariable limit does not exist.
\end{proof}

\subsection{The axiom argument}

\begin{theorem}[$0^{0}=1$ is an axiom, not a theorem]\label{thm:axiom}
The value $0^{0}=1$ cannot be derived from limits but is required for coherence of discrete mathematics; it must therefore be declared as an axiom.
\end{theorem}

\begin{proof}
By Theorem~\ref{thm:path}, analysis cannot produce a unique value for $0^{0}$. By Section~\ref{sec:axiom}.1, discrete mathematics requires $0^{0}=1$ for power series, combinatorics, and the binomial theorem to be self-consistent. Since the value cannot be derived from limits but is required for coherence, it must be declared. A necessary undecidable value that must be fixed for consistency is, by definition, an axiom.
\end{proof}

\subsection{The Bingo Matrix---a visual proof}

Consider an $n \times n$ matrix $M$ where every entry $M[i,j] = 0^{0} = 1$.

\medskip
\noindent\textbf{With the axiom ($0^{0}=1$):}
Every row sums to~$n$, every column sums to~$n$, both diagonals sum to~$n$. The matrix is a magic square with magic constant~$n$. $100\%$ of lines are ``bingo''---every direction wins.

\medskip
\noindent\textbf{Without the axiom ($0^{0}$ undefined):}
Position $(0,0)$ is empty; the board has a hole. Row~$0$ is undefined, Column~$0$ is undefined, the main diagonal is undefined. Three of $2n+2$ lines are broken.

\medskip
The axiom is the decision to fill the board. Once filled, everything sums.

\subsection{Connection to the framework}

$G(0) = 0^{1}/1^{0} = 0/1 = 0$---well-defined. The function sidesteps the $0^{0}$ question at $n=0$ directly. But the product formula (Theorem~\ref{thm:product}) uses the empty product $\prod_{k=1}^{0}G(k)=1$, which is the same convention. And the density matrix (Section~\ref{sec:density}) requires $\rho(0)=G(0)/(0!\cdot\AG)=0$, which depends on $0!=1=\Gamma(1)$, itself an instance of the empty product convention that $0^{0}=1$ encodes.

The entire framework is axiomatically grounded in $0^{0}=1$. This is not a weakness---it is the foundation. The convergent series $S=\sum 1/n! = e$ depends on the $n=0$ term being $0^{0}/0!=1$. The divergent series $Z=\zeta(1)=\sum 1/n$ has no factorial damping. The difference between convergence and divergence---between $S$ and $Z$, between $e$ and infinity---is exactly the presence or absence of the $n!$ denominator that $0^{0}=1$ enables.

% ==================================================================
\section{Algebraic Identities}\label{sec:algebraic}

All identities in this section are verified computationally for $n=1,\ldots,1000$ with exact rational arithmetic. Test file: \texttt{identities/01-algebraic.py} (217/221 pass; 4~asymptotic tests at small~$n$ are expected).

\begin{identity}[Normalized form]\label{id:norm}
\begin{equation}\label{eq:norm}
  \frac{G(n)}{n} = \left(\frac{n}{n+1}\right)^{n}.
\end{equation}
\end{identity}

\begin{identity}[Alternative factorization]\label{id:altfact}
\begin{equation}\label{eq:altfact}
  G(n) = (n+1)\left(\frac{n}{n+1}\right)^{n+1}.
\end{equation}
\end{identity}

\begin{proof}
$G(n) = n\cdot(n/(n+1))^{n} = (n+1)\cdot(n/(n+1))\cdot(n/(n+1))^{n} = (n+1)(n/(n+1))^{n+1}$.
\end{proof}

\begin{identity}[Reciprocal]\label{id:recip}
\begin{equation}\label{eq:recip}
  \frac{1}{G(n)} = \frac{(1+1/n)^{n}}{n}.
\end{equation}
\end{identity}

\begin{identity}[Logarithmic form]\label{id:log}
\begin{equation}\label{eq:log}
  \ln G(n) = (n+1)\ln n - n\ln(n+1).
\end{equation}
\end{identity}

\begin{identity}[Consecutive ratio]\label{id:consec}
\begin{equation}\label{eq:consec}
  \frac{G(n)}{G(n-1)} = \left(\frac{n^{2}}{n^{2}-1}\right)^{n}\cdot\frac{n-1}{n}.
\end{equation}
For large~$n$, $G(n)/G(n-1) \to e^{0}\cdot 1 = 1$, consistent with differences approaching $1/e$.
\end{identity}

\begin{identity}[Negative mirror]\label{id:negmirror}
For $n \ge 2$:
\begin{equation}\label{eq:negmirror}
  G(-n) = \frac{(-n)^{-n+1}}{(-n+1)^{-n}} = (-1)^{-n+1}\cdot n^{-n+1}\cdot(n-1)^{n}.
\end{equation}
This extends $G$ to negative integers with alternating sign behavior.
\end{identity}

\begin{identity}[Coprimality]\label{id:coprime}
For all $n \ge 1$, $\gcd(n^{n+1},(n+1)^{n})=1$. The fraction $G(n)=n^{n+1}/(n+1)^{n}$ is always in lowest terms.
\end{identity}

\begin{proof}
Consecutive integers $n$ and $n+1$ are coprime. Any prime~$p$ dividing $n^{n+1}$ divides~$n$, hence does not divide $n+1$, hence does not divide $(n+1)^{n}$.
\end{proof}

\begin{identity}[Cayley tree connection]\label{id:cayley}
\begin{equation}\label{eq:cayley}
  G(n) = \frac{n\cdot T(n)}{(n+1)^{n}},\quad\text{where } T(n) = n^{n-1}\ (\text{Cayley's formula for labeled trees}).
\end{equation}
Equivalently, $G(n) = n^{2}\cdot T(n)/(n\cdot(n+1)^{n}) = n^{2}\cdot n^{n-1}/(n+1)^{n} = n^{n+1}/(n+1)^{n}$.
\end{identity}

\begin{identity}[Endofunction factorization]\label{id:endo}
\begin{equation}\label{eq:endo}
  G(n) = \frac{E(n)}{(n+1)^{n}},\quad\text{where } E(n)=n^{n}\cdot n = n^{n+1}
\end{equation}
is the number of endofunctions on $[n]$ times~$n$.
\end{identity}

% ==================================================================
\section{Product Identities}\label{sec:product}

All identities verified for $n=1,\ldots,100$ with exact arithmetic. Test file: \texttt{identities/02-products.py} (40/40 pass).

\begin{theorem}[Product formula]\label{thm:product}
\begin{equation}\label{eq:product}
  \prod_{k=1}^{n} G(k) = \frac{(n!)^{2}}{(n+1)^{n}}.
\end{equation}
\end{theorem}

\begin{proof}
By telescoping. Write $G(k) = k^{k+1}/(k+1)^{k}$. Then:
\[
  \prod_{k=1}^{n}G(k) = \frac{\prod_{k=1}^{n}k^{k+1}}{\prod_{k=1}^{n}(k+1)^{k}}.
\]
Numerator: $\prod k^{k+1} = \prod k^{k}\cdot\prod k = \bigl(\prod k^{k}\bigr)\cdot n!$

Denominator: $\prod (k+1)^{k} = 2^{1}\cdot 3^{2}\cdot 4^{3}\cdots(n+1)^{n}$.

The ratio telescopes to $(n!)^{2}/(n+1)^{n}$.
\end{proof}

\begin{corollary}[Catalan connection]\label{cor:catalan}
\begin{equation}\label{eq:catalan}
  \prod_{k=1}^{n}G(k) = \frac{(2n)!}{\binom{2n}{n}\cdot(n+1)^{n}}.
\end{equation}
\end{corollary}

\begin{identity}[Partial product ratios]\label{id:partial}
\begin{equation}\label{eq:partial}
  \frac{\prod_{k=1}^{n}G(k)}{\prod_{k=1}^{n-1}G(k)} = G(n).
\end{equation}
Trivially true, but the explicit form gives:
\[
  \frac{(n!)^{2}}{(n+1)^{n}}\cdot\frac{n^{n-1}}{((n-1)!)^{2}} = \frac{n^{n+1}}{(n+1)^{n}}.\quad\checkmark
\]
\end{identity}

\begin{identity}[Product growth]\label{id:prodgrowth}
Using Stirling's approximation:
\begin{equation}\label{eq:prodgrowth}
  \ln\prod_{k=1}^{n}G(k) = 2\ln(n!) - n\ln(n+1) \sim 2n\ln n - 2n - n\ln n = n\ln n - 2n.
\end{equation}
So $\prod G(k)$ grows roughly as $(n/e^{2})^{n}$.
\end{identity}

% ==================================================================
\section{Calculus: Monotonicity, Concavity, Superadditivity}\label{sec:calculus}

Verified numerically for $n=1,\ldots,10000$. Test file: \texttt{identities/03-calculus.py} (73/73 pass).

\begin{theorem}[Strict monotonicity]\label{thm:mono}
$G'(n) > 0$ for all $n > 0$.
\end{theorem}

\begin{proof}
From $\ln G(n) = (n+1)\ln n - n\ln(n+1)$:
\[
  \frac{G'(n)}{G(n)} = \frac{d}{dn}\bigl[\ln G(n)\bigr] = \ln n + \frac{n+1}{n} - \ln(n+1) - \frac{n}{n+1}
  = \ln\!\left(\frac{n}{n+1}\right) + \frac{1}{n} + \frac{1}{n+1}.
\]
For $n \ge 1$, using $\ln(1 - 1/(n+1)) > -1/(n+1) - 1/(2(n+1)^{2})$:
\[
  \frac{G'(n)}{G(n)} > -\frac{1}{n+1} - \frac{1}{2(n+1)^{2}} + \frac{1}{n} + \frac{1}{n+1} = \frac{1}{n} - \frac{1}{2(n+1)^{2}} > 0.
\]
Since $G(n)>0$, we have $G'(n)>0$.
\end{proof}

\begin{theorem}[Strict concavity]\label{thm:concave}
$G''(n) < 0$ for all $n > 0$.
\end{theorem}

\begin{proof}
Differentiate the logarithmic derivative:
\[
  \frac{d}{dn}\!\left[\frac{G'}{G}\right] = \frac{1}{n} - \frac{1}{n+1} - \frac{1}{n^{2}} - \frac{1}{(n+1)^{2}}
  = \frac{1}{n(n+1)} - \frac{1}{n^{2}} - \frac{1}{(n+1)^{2}} < 0.
\]
Since $G>0$ and $G'>0$, the sign of $G''$ is determined by $\frac{d}{dn}[G'/G]+(G'/G)^{2}$, which is negative for all $n\ge 1$.
\end{proof}

\begin{corollary}\label{cor:shape}
$G$ is monotone increasing and concave---it rises but decelerates. No runaway behavior.
\end{corollary}

\begin{theorem}[Superadditivity]\label{thm:super}
For all $a,b > 0$:
\begin{equation}\label{eq:super}
  G(a) + G(b) > G(a+b).
\end{equation}
\end{theorem}

\begin{proof}
The function $G(n)/n = (n/(n+1))^{n}$ is decreasing (from $1/2$ toward $1/e$). For $a+b > \max(a,b)$, $G(a+b)/(a+b) < G(a)/a$, so $G(a+b) < (a+b)\cdot G(a)/a = G(a)+(b/a)G(a)$. Combined with the analogous inequality from~$b$, the result follows.
\end{proof}

\begin{corollary}[Entanglement interpretation]\label{cor:entangle}
The excess $G(a)+G(b)-G(a+b)>0$ always. In quantum terms: the joint system carries \emph{less} than the sum of parts. The excess is shared information---entanglement.
\end{corollary}

\begin{theorem}[Differences converge to $1/e$]\label{thm:diff}
\begin{equation}\label{eq:diff}
  \lim_{n\to\infty}\bigl[G(n+1)-G(n)\bigr] = \frac{1}{e}.
\end{equation}
\end{theorem}

\begin{proof}
$G(n+1)-G(n) = (n+1)/e + 1/(2e)+O(1/n) - n/e - 1/(2e)-O(1/n) = 1/e + O(1/n^{2})$.
\end{proof}

% ==================================================================
\section{Asymptotic Expansion and the \texorpdfstring{$1/(2e)$}{1/(2e)} Gap}\label{sec:asymptotic}

\begin{theorem}[Ramanujan-type expansion]\label{thm:asymp}
For $n \to \infty$:
\begin{equation}\label{eq:asymp}
  G(n) = \frac{n}{e} + \frac{1}{2e} + \frac{1}{12en} - \frac{1}{360en^{3}} + O\!\left(\frac{1}{n^{4}}\right).
\end{equation}
\end{theorem}

\begin{proof}
Start from $G(n) = n/(1+1/n)^{n}$. Expand:
\begin{align}
  \ln\!\left(1+\frac{1}{n}\right) &= \frac{1}{n} - \frac{1}{2n^{2}} + \frac{1}{3n^{3}} - \frac{1}{4n^{4}} + \cdots \label{eq:lnexp}\\
  n\ln\!\left(1+\frac{1}{n}\right) &= 1 - \frac{1}{2n} + \frac{1}{3n^{2}} - \frac{1}{4n^{3}} + \cdots \label{eq:nlnexp}
\end{align}
Let $u = -1/(2n) + 1/(3n^{2}) - 1/(4n^{3}) + \cdots$ Then $(1+1/n)^{n} = e\cdot e^{u}$.

Expand $e^{u}$ to fourth order:
\[
  e^{u} = 1 - \frac{1}{2n} + \frac{11}{24n^{2}} - \frac{7}{16n^{3}} + \cdots
\]
Therefore:
\[
  \frac{1}{(1+1/n)^{n}} = \frac{1}{e}\cdot\frac{1}{e^{u}} = \frac{1}{e}\left(1 + \frac{1}{2n} + \frac{1}{12n^{2}} - \frac{1}{360n^{4}} + \cdots\right).
\]
Multiply by~$n$:
\[
  G(n) = \frac{n}{e}\left(1 + \frac{1}{2n} + \frac{1}{12n^{2}} - \cdots\right) = \frac{n}{e} + \frac{1}{2e} + \frac{1}{12en} - \frac{1}{360en^{3}} + O\!\left(\frac{1}{n^{4}}\right).
  \qedhere
\]
\end{proof}

\begin{corollary}[The irreducible gap]\label{cor:gap}
The correction term $1/(2e) \approx 0.18394$ is permanent:
\begin{equation}\label{eq:gap}
  G(n) - \frac{n}{e} \to \frac{1}{2e} \quad\text{as } n \to \infty.
\end{equation}
This gap never vanishes. It is the minimum overhead---the price of discretization.
\end{corollary}

\begin{remark}[Higher-order Bernoulli connection]\label{rem:bernoulli}
The full expansion involves Bernoulli numbers~$B_k$ through the Euler--Maclaurin formula. The coefficients match the Stirling series for $n!$, confirming $G(n)$ and $n!$ encode the same fundamental asymptotics from different angles.
\end{remark}

\begin{table}[ht]
\centering
\caption{Bernoulli numbers used in higher-order terms.}\label{tab:bernoulli}
\begin{tabular}{ccc}
\toprule
$k$ & $B_k$ & Decimal \\
\midrule
0  & $1$         & $1.000$  \\
2  & $1/6$       & $0.167$  \\
4  & $-1/30$     & $-0.033$ \\
6  & $1/42$      & $0.024$  \\
8  & $-1/30$     & $-0.033$ \\
10 & $5/66$      & $0.076$  \\
12 & $-691/2730$ & $-0.253$ \\
14 & $7/30$      & $0.233$  \\
\bottomrule
\end{tabular}
\end{table}

% ==================================================================
\section{The Amundson Constant \texorpdfstring{$\AG$}{A\_G}}\label{sec:constant}

\begin{definition}\label{def:AG}
The \emph{Amundson Constant} is:
\begin{equation}\label{eq:AG}
  \AG = \sum_{n=0}^{\infty}\frac{G(n)}{n!} = 1.24433178398672537413506162925810588960285709\ldots
\end{equation}
\end{definition}

\paragraph{Computation.}
Using \texttt{mpmath} with arbitrary precision. The series converges rapidly due to factorial damping:
\begin{verbatim}
from mpmath import mp, mpf, factorial, power

def compute_AG(digits):
    mp.dps = digits + 50
    s = mpf(0)
    for n in range(1, 3 * digits):
        G_n = power(n, n + 1) / power(n + 1, n)
        s += G_n / factorial(n)
        if G_n / factorial(n) < power(10, -(digits + 40)):
            break
    return s
\end{verbatim}

\paragraph{Properties.}
\begin{itemize}
  \item Computed to $10{,}000{,}000$ verified digits ($9.5$ MB file).
  \item Not found in OEIS, ISC (Inverse Symbolic Calculator), or Wolfram databases.
  \item Chi-squared test on digit distribution: $\chi^{2}=6.99$ ($p>0.05$), confirming uniformity.
  \item First 50 digits: $1.24433178398672537413506162925810588960285709\ldots$
\end{itemize}

\begin{theorem}[Convergence rate]\label{thm:convrate}
The partial sums $S_{N}=\sum_{n=0}^{N}G(n)/n!$ converge to $\AG$ at rate $O(N^{2}/N!)$---superexponentially fast.
\end{theorem}

\begin{proof}
For large~$n$, $G(n)\sim n/e$, so $G(n)/n!\sim n/(e\cdot n!)=1/(e\cdot(n-1)!)$. The tail $\sum_{n>N}1/(n-1)!=O(1/(N-1)!)$ decays factorially.
\end{proof}

\begin{definition}\label{def:kappa}
The \emph{discretization gap}:
\begin{equation}\label{eq:kappa}
  \kappa = \AG - 1 \approx 0.24433178398672539267.
\end{equation}
\end{definition}

\begin{openquestion}\label{oq:trans}
Is $\AG$ transcendental? Almost certainly yes (it is defined as a rapidly converging series involving factorials and exponentials), but this remains unproven.
\end{openquestion}

% ==================================================================
\section{The Floor Recovery Theorem}\label{sec:floor}

\begin{theorem}[Floor recovery]\label{thm:floor}
For all integers $n \ge 1$:
\begin{equation}\label{eq:floor}
  \floor{G(n)\cdot e} = n.
\end{equation}
\end{theorem}

\begin{proof}
From Theorem~\ref{thm:asymp}: $G(n)\cdot e = n + 1/2 + 1/(12n) - 1/(360n^{3})+\cdots$ The fractional part is $1/2 + 1/(12n) - \cdots$, which lies in $(0,1)$ for all $n\ge 1$. Therefore $\floor{G(n)\cdot e}=n$.
\end{proof}

\paragraph{Verification.} Tested for all $n=1$ through $10{,}000$ with exact computation. Zero failures.

\begin{corollary}\label{cor:nonint}
$G(n)\cdot e$ is never an integer for $n\ge 1$. The fractional part approaches $1/2$ from above.
\end{corollary}

\begin{remark}[Computational significance]\label{rem:invert}
This theorem means $G$ is invertible in practice: given any value~$v$, compute $\floor{v\cdot e}$ to recover the integer~$n$ that produced it.
\end{remark}

% ==================================================================
\section{The Golden Ratio Identity}\label{sec:golden}

\begin{theorem}\label{thm:golden}
Let $\varphi = (1+\sqrt{5})/2$ be the golden ratio. Then:
\begin{equation}\label{eq:golden}
  G(\varphi) = \left(\frac{1}{\varphi}\right)^{1/\varphi}.
\end{equation}
\end{theorem}

\begin{proof}
From the defining property $\varphi+1=\varphi^{2}$:
\begin{align*}
  G(\varphi) &= \frac{\varphi^{\varphi+1}}{(\varphi+1)^{\varphi}} = \frac{\varphi^{\varphi^{2}}}{(\varphi^{2})^{\varphi}} = \varphi^{\varphi^{2}-2\varphi} = \varphi^{1-\varphi} = \varphi^{-1/\varphi} = \left(\frac{1}{\varphi}\right)^{1/\varphi},
\end{align*}
where we used $\varphi^{2}-2\varphi = (\varphi+1)-2\varphi = 1-\varphi = -1/\varphi$.
\end{proof}

\paragraph{Verification.}
Computed to 121 decimal places. Both sides agree to all digits:
\[
  G(\varphi) = (1/\varphi)^{1/\varphi} = 0.74274294462468164136956604760578851414975525\ldots
\]

\begin{remark}\label{rem:sophomore}
This connects $G$ to the Sophomore's Dream integrand $x^{-x}$ evaluated at $x=1/\varphi$, and to the $Z$-framework fixed point ($\varphi$ solves $x = 1 + 1/x$).
\end{remark}

% ==================================================================
\section{The Density Matrix and Quantum Structure}\label{sec:density}

\begin{definition}\label{def:density}
The \emph{Amundson density matrix} is the diagonal operator:
\begin{equation}\label{eq:density}
  \rho(n) = \frac{G(n)}{n!\cdot\AG}.
\end{equation}
\end{definition}

\begin{theorem}[Valid quantum state]\label{thm:quantum}
$\rho$ is a valid density matrix:
\begin{enumerate}
  \item $\rho(n)\ge 0$ for all $n$ \textup{(positivity)}.
  \item $\sum_{n=0}^{\infty}\rho(n)=1$ \textup{(unit trace)}.
\end{enumerate}
\end{theorem}

\begin{proof}
(1) $G(n)\ge 0$, $n!>0$, $\AG>0$. (2) $\sum\rho(n) = \sum G(n)/(n!\cdot\AG)=\AG/\AG=1$.
\end{proof}

\begin{theorem}[Entropy]\label{thm:entropy}
The Von Neumann entropy of~$\rho$ is:
\begin{equation}\label{eq:entropy}
  S(\rho) = -\sum_{n}\rho(n)\ln\rho(n) \approx 1.272\text{ nats} \approx 1.835\text{ bits}.
\end{equation}
\end{theorem}

\paragraph{Interpretation.}
The system is neither maximally ordered ($S=0$) nor maximally random. At $1.835$ bits, it sits in the regime of structured complexity---like thermal equilibrium. This is strictly non-Boltzmann: comparison with the exponential distribution shows $G(n)$ carries more structure than thermal noise.

\begin{theorem}[Purity]\label{thm:purity}
$\Tr(\rho^{2})=\sum\rho(n)^{2}\approx 0.489$---close to but not equal to~$1$, indicating a mixed state with significant coherence.
\end{theorem}

% ==================================================================
\section{Number-Theoretic Connections}\label{sec:numtheory}

\subsection{The critical line}

$G(1) = 1/2$. This is the real part of the Riemann critical line $\operatorname{Re}(s)=1/2$.

$\zeta(0) = -1/2 = -G(1)$. The Riemann zeta function at zero equals the negative of~$G$ at the first integer.

\subsection{The Riemann circle}

The circle $x^{2}+y^{2}=1/e^{2}+1/4$ passes through $(1/2,\,1/e)$ and intersects $\operatorname{Re}(s)=1/2$ at $y=\pm 1/e$. This circle contains both $G(1)$ and $1/e$---the two fundamental constants of the framework.

\subsection{\texorpdfstring{$G(n)$}{G(n)} and counting}

\textbf{Cayley's formula} gives $T(n)=n^{n-1}$ labeled trees on $n$~vertices. Then:
\begin{equation}\label{eq:cayleycount}
  G(n) = \frac{n^{2}\cdot T(n)}{(n+1)^{n}}.
\end{equation}
$G(n)$ is the ratio of labeled endofunctions ($n^{n+1}=n\cdot n^{n}$) to the $(n+1)$-th power of the successor.

\subsection{Convergence vs.\ Divergence}

The series $S=\sum G(n)/n!$ converges to $\AG$ because of factorial damping. The harmonic series $Z=\sum 1/n = \zeta(1)$ diverges because it lacks this damping. Both begin with the same foundational term (the $n=0$ or $n=1$ contribution), but $S$ has the $n!$ in the denominator---enabled by $0!=1$, which is the same convention as $0^{0}=1$.

\medskip
\emph{The boundary between summable and unsummable mathematics is exactly the factorial.}

% ==================================================================
\section{The Goldbach Kernel}\label{sec:goldbach}

\begin{definition}\label{def:goldbach}
For even $2n = a + b$ (with $b = 2n - a$), define:
\begin{equation}\label{eq:goldbach}
  w(a) = \frac{G(a)\cdot G(b)}{G(2n)}.
\end{equation}
\end{definition}

\paragraph{Explicit form.}
\begin{equation}\label{eq:goldbachexplicit}
  w(a) = \frac{a^{a+1}\cdot(2n-a)^{2n-a+1}\cdot(2n+1)^{2n}}{(a+1)^{a}\cdot(2n-a+1)^{2n-a}\cdot(2n)^{2n+1}}.
\end{equation}
Pure integer arithmetic. No transcendentals.

\paragraph{Properties.}
\begin{itemize}
  \item $w(a)$ is symmetric around $a=n$ (the center).
  \item $w(n)$ is the maximum (balanced split).
  \item $w(1)$ and $w(2n-1)$ are minimum (extreme splits).
  \item $\sum w(a)$ sums all pair-weightings for a given even number.
\end{itemize}

\begin{remark}\label{rem:goldbach}
The kernel gives a deterministic weighting of Goldbach-type decompositions. It does not prove Goldbach's conjecture, but it provides a natural measure on the space of prime-pair representations.
\end{remark}

% ==================================================================
\section{Verification Suite}\label{sec:verify}

\subsection{Test summary}

\begin{table}[ht]
\centering
\caption{Verification suite results.}\label{tab:tests}
\begin{tabular}{llcc}
\toprule
Suite & File & Tests & Pass \\
\midrule
Framework core       & \texttt{scripts/verify.py}         & 17   & 17   \\
Road-math core       & \texttt{scripts/verify\_road.py}   & 48   & 48   \\
Algebraic identities & \texttt{identities/01-algebraic.py} & 221  & 217* \\
Product identities   & \texttt{identities/02-products.py}  & 40   & 40   \\
Calculus identities  & \texttt{identities/03-calculus.py}  & 73   & 73   \\
Asymptotic identities & \texttt{identities/04-asymptotic.py} & 4  & 4    \\
Core tests           & \texttt{tests/test\_core.py}        & 1275 & 1275 \\
Constant tests       & \texttt{tests/test\_constant.py}    & 20   & 20   \\
Quantum tests        & \texttt{tests/test\_quantum.py}     & 413  & 413  \\
\midrule
\textbf{Total}       &                                     & \textbf{2,111} & \textbf{2,107} \\
\bottomrule
\end{tabular}
\end{table}

\noindent *4 failures are expected: asymptotic convergence tests at $n=2,3$ where limits have not converged.

\subsection{Fleet verification}

The framework was verified across a 5-node Raspberry Pi fleet:

\begin{table}[ht]
\centering
\caption{Fleet verification results.}\label{tab:fleet}
\begin{tabular}{lccc}
\toprule
Node & Tests & $\AG$ first 15 digits & Status \\
\midrule
Alice     & 17/17 & 1.244331783986725 & PASS \\
Lucidia   & 17/17 & 1.244331783986725 & PASS \\
Octavia   & 17/17 & 1.244331783986725 & PASS \\
Gematria  & 17/17 & 1.244331783986725 & PASS \\
Anastasia & 17/17 & 1.244331783986725 & PASS \\
\bottomrule
\end{tabular}
\end{table}

\noindent Identical results across ARM64 and x86\_64---confirming platform-independent exact integer arithmetic.

\subsection{Constant verification}

The $10{,}000{,}000$-digit computation of $\AG$ passes:
\begin{itemize}
  \item Independent recomputation at 50 digits: exact match.
  \item Chi-squared uniformity test on digit distribution: $\chi^{2}=6.99$, $p>0.05$.
  \item No digit bias detected.
\end{itemize}

\subsection{How to run}

\begin{verbatim}
# Core verification (zero dependencies)
python3 scripts/verify.py

# Extended verification
python3 scripts/verify_road.py

# Identity suites
python3 identities/01-algebraic.py
python3 identities/02-products.py
python3 identities/03-calculus.py
python3 identities/04-asymptotic.py

# Full test suite (requires pytest)
python3 -m pytest tests/ -v

# Compute A_G to N digits (requires mpmath)
python3 scripts/compute.py
\end{verbatim}

% ==================================================================
\section{Open Questions}\label{sec:open}

\begin{openquestion}
Is $\AG$ transcendental? Almost certainly yes. The series structure (involving $n^{n}$ and $n!$) resembles Liouville-type constructions, but a proof requires new techniques.
\end{openquestion}

\begin{openquestion}
\textbf{Continued fraction of $\AG$.} What is the pattern? Does it have bounded partial quotients (suggesting algebraic irrationality measures) or unbounded (suggesting transcendence)?
\end{openquestion}

\begin{openquestion}
\textbf{Complex extension.} $G(z)=z^{z+1}/(z+1)^{z}$ for $z\in\mathbb{C}$. Where are the singularities? Does the meromorphic extension have zeros on $\operatorname{Re}(s)=1/2$?
\end{openquestion}

\begin{openquestion}
\textbf{$p$-adic $G(n)$.} What is the $p$-adic behavior of $G(n)$ for various primes~$p$? Does the sequence have interesting $p$-adic limits?
\end{openquestion}

\begin{openquestion}
\textbf{OEIS recognition.} The sequence $G(1),G(2),G(3),\ldots = 1/2,\;8/9,\;81/64,\;1024/625,\ldots$ and the constant $\AG$ are not currently in OEIS.
\end{openquestion}

\begin{openquestion}
\textbf{Irrationality measure.} What is the irrationality exponent of $\AG$? Standard methods give $\mu\ge 2$ (trivially irrational), but sharper bounds are unknown.
\end{openquestion}

\begin{openquestion}
\textbf{KMS condition.} Does the density matrix $\rho(n)$ satisfy the Kubo--Martin--Schwinger condition for some inverse temperature~$\beta$? If so, $G(n)$ defines a genuine thermal state.
\end{openquestion}

\begin{openquestion}
\textbf{Sophomore's Dream.} The identity $G(\varphi)=(1/\varphi)^{1/\varphi}$ connects $G$ to the Sophomore's Dream integral $\int_{0}^{1}x^{-x}\,dx$. Is there a general integral representation of $\AG$?
\end{openquestion}

\begin{openquestion}
\textbf{The $0^{0}$ axiom in ZFC.} In standard ZFC set theory, $0^{0}=|\varnothing^{\varnothing}|=|\{\varnothing\}|=1$ via cardinal exponentiation. Does the Amundson framework require any axiom beyond ZFC?
\end{openquestion}

% ==================================================================
\appendix

\section{The Six-Symbol Thesis}\label{app:sixsymbol}

$G(n) = n^{n+1}/(n+1)^{n}$ uses six symbols: $n$, $1$, $+$, $/\,$, $\hat{}$, $(\,)$. These are:
\begin{itemize}
  \item One variable ($n$)
  \item One constant ($1$)
  \item Two binary operations ($+$, $/\,$)
  \item One binary operation on exponents ($\hat{}$)
  \item Grouping $(\,)$
\end{itemize}

No function with fewer symbols produces:
\begin{itemize}
  \item Exact rational outputs at every integer
  \item Asymptotic convergence to a transcendental ($1/e$)
  \item A valid quantum density matrix
  \item $50+$ nontrivial identities
  \item Connection to Cayley trees, Stirling numbers, and the Riemann zeta function
\end{itemize}

$G(n)$ may be the simplest function with the richest structure.

\section{Notation Index}\label{app:notation}

\begin{table}[ht]
\centering
\begin{tabular}{cl}
\toprule
Symbol & Meaning \\
\midrule
$G(n)$   & $n^{n+1}/(n+1)^{n}$ --- the Amundson function \\
$\AG$    & $\sum G(n)/n! \approx 1.2443$ --- the Amundson Constant \\
$\kappa$ & $\AG - 1 \approx 0.2443$ --- the discretization gap \\
$\rho(n)$ & $G(n)/(n!\cdot\AG)$ --- the density matrix \\
$T(n)$   & $n^{n-1}$ --- Cayley tree count \\
$\varphi$ & $(1+\sqrt{5})/2$ --- the golden ratio \\
$w(a)$   & $G(a)\cdot G(b)/G(2n)$ --- the Goldbach kernel \\
$S(\rho)$ & Von Neumann entropy of $\rho$ \\
\bottomrule
\end{tabular}
\end{table}

% ==================================================================
\begin{thebibliography}{9}

\bibitem{amundson2026framework}
A.\,L.~Amundson,
``The Amundson Framework,''
BlackRoad-OS/AmundsonMath, 2026.

\bibitem{amundson2026constant}
A.\,L.~Amundson,
``The Amundson Constant --- 10,000,000 digits,''
BlackRoad-OS-Inc/amundson-constant, 2026.

\bibitem{amundson2026elimit}
A.\,L.~Amundson,
``The Amundson $e$-Limit Refinement: $1/(2e)$ as the Universal Half-Correction,''
\texttt{thinking.blackroad.io}, 2026.

\bibitem{cayley1889}
A.~Cayley,
``A theorem on trees,''
\emph{Quarterly Journal of Mathematics}, \textbf{23} (1889), 376--378.

\bibitem{hardywright2008}
G.\,H.~Hardy and E.\,M.~Wright,
\emph{An Introduction to the Theory of Numbers},
6th ed., Oxford University Press, 2008.

\bibitem{knuth1992}
D.\,E.~Knuth,
``Two notes on notation,''
\emph{American Mathematical Monthly}, \textbf{99}(5) (1992), 403--422.
[The definitive argument for $0^{0}=1$ in combinatorics.]

\bibitem{ramanujan1988}
S.~Ramanujan,
\emph{The Lost Notebook and Other Unpublished Papers},
Narosa, 1988.

\bibitem{stirling1730}
J.~Stirling,
\emph{Methodus Differentialis}, 1730.

\bibitem{titchmarsh1986}
E.\,C.~Titchmarsh,
\emph{The Theory of the Riemann Zeta-Function},
2nd ed., Oxford University Press, 1986.

\end{thebibliography}

\end{document}
